\documentclass[11pt, a4paper]{article}

% --- PAQUETES PRINCIPALES ---
\usepackage[utf8]{inputenc}
\usepackage[spanish, es-tabla]{babel}
\usepackage[margin=2.5cm]{geometry}
\usepackage{amsmath, amssymb, mathtools}
\usepackage{siunitx}
\usepackage{graphicx}
\usepackage{booktabs}
\usepackage{tabularx}
\usepackage[most]{tcolorbox}
\usepackage{listings}
\usepackage{xcolor}

% --- PAQUETES PARA GRÁFICAS Y CIRCUITOS ---
\usepackage{pgfplots}
\usepackage[american]{circuitikz}
\pgfplotsset{compat=1.18}

% --- CONFIGURACIÓN DE CÓDIGO (PYTHON) ---
\definecolor{codegreen}{rgb}{0,0.6,0}
\definecolor{codegray}{rgb}{0.5,0.5,0.5}
\definecolor{codepurple}{rgb}{0.58,0,0.82}
\definecolor{backcolour}{rgb}{0.96,0.96,0.96}
\lstset{
    backgroundcolor=\color{backcolour},   
    commentstyle=\color{codegreen},
    keywordstyle=\color{magenta},
    numberstyle=\tiny\color{codegray},
    stringstyle=\color{codepurple},
    basicstyle=\ttfamily\footnotesize,
    breaklines=true,                 
    numbers=left,                    
    numbersep=5pt,                  
    language=Python,
    literate={á}{{\'a}}1 {é}{{\'e}}1 {í}{{\'i}}1 {ó}{{\'o}}1 {ú}{{\'u}}1
             {Á}{{\'A}}1 {É}{{\'E}}1 {Í}{{\'I}}1 {Ó}{{\'O}}1 {Ú}{{\'U}}1
             {ñ}{{\~n}}1 {Ñ}{{\~N}}1
}

% --- CONFIGURACIÓN DE SIUNITX ---
\sisetup{output-decimal-marker = {,}, inter-unit-product = \ensuremath{{}\cdot{}}}

% --- CABECERAS ---
\usepackage{fancyhdr}
\pagestyle{fancy}
\fancyhf{}
\lhead{\textbf{Circuitos Electrónicos III (IMT-221)}}
\rhead{Sesión 4: Modelos y Pequeña Señal}
\cfoot{\thepage}
\renewcommand{\headrulewidth}{0.4pt}

% --- ENTORNOS PERSONALIZADOS ---
\newtcolorbox{aviso}[1][]{colback=red!5!white, colframe=red!75!black, fonttitle=\bfseries, title=#1}

\begin{document}

\begin{center}
    {\LARGE \textbf{Apuntes de Cátedra: Sesión 4}}\\[0.3cm]
    {\large \textbf{Jerarquía de Modelos, Pequeña Señal en CA y Límite de Linealidad}}\\
    \textit{Docente: M.Sc. Ing. Bernardo Quiroga Turdera \quad | \quad Fecha: 14 de agosto de 2026}
    \vspace{0.3cm}
    \hrule
\end{center}

\vspace{0.3cm}

% ==========================================
% SECCIÓN 1: JERARQUÍA DE MODELOS EN DC
% ==========================================
\section{Jerarquía de Modelos del Diodo en DC (Gran Señal)}
Para resolver circuitos no lineales sin recurrir a iteraciones complejas, la ingeniería utiliza tres modelos equivalentes que aproximan la curva exponencial de Shockley, dependiendo de la precisión requerida.

\vspace{0.3cm}

\noindent
\begin{minipage}{0.32\textwidth}
    \centering \textbf{1. Modelo Ideal}
    \vspace{0.2cm}\\
    \begin{circuitikz}[scale=0.7, transform shape]
        \draw (0,0) to[Do, l=Ideal] (3,0);
        \draw[fill=white, draw=none] (0,-1) rectangle (3,1);
    \end{circuitikz}\\
    \vspace{0.2cm}
    \begin{tikzpicture}[scale=0.6]
        \draw[->] (-1.5,0) -- (1.5,0) node[right] {$V_D$};
        \draw[->] (0,-0.5) -- (0,2.5) node[above] {$I_D$};
        \draw[ultra thick, blue] (-1.5,0) -- (0,0) -- (0,2);
    \end{tikzpicture}\\
    \vspace{0.2cm}
    \raggedright\footnotesize
    \textbf{Directa:} Interruptor cerrado ($V_D=0$).\\
    \textbf{Inversa:} Abierto ($I_D=0$).\\
    \textbf{Uso:} Análisis cualitativo rápido y compuertas lógicas simples.
\end{minipage}
\hfill
\begin{minipage}{0.32\textwidth}
    \centering \textbf{2. Caída Constante ($V_\gamma$)}
    \vspace{0.2cm}\\
    \begin{circuitikz}[scale=0.7, transform shape]
        \draw (0,0) to[Do] (1.5,0) to[V, v=$0.7\text{V}$] (3,0);
    \end{circuitikz}\\
    \vspace{0.2cm}
    \begin{tikzpicture}[scale=0.6]
        \draw[->] (-0.5,0) -- (2,0) node[right] {$V_D$};
        \draw[->] (0,-0.5) -- (0,2.5) node[above] {$I_D$};
        \draw[ultra thick, blue] (-0.5,0) -- (1,0) -- (1,2);
        \node[below] at (1,0) {$0.7\text{V}$};
    \end{tikzpicture}\\
    \vspace{0.2cm}
    \raggedright\footnotesize
    \textbf{Directa:} Batería fija de \qty{0.7}{\volt}.\\
    \textbf{Inversa:} Abierto.\\
    \textbf{Uso:} Diseño estándar de fuentes de alimentación y rectificadores.
\end{minipage}
\hfill
\begin{minipage}{0.32\textwidth}
    \centering \textbf{3. Lineal por Tramos (PWL)}
    \vspace{0.2cm}\\
    \begin{circuitikz}[scale=0.7, transform shape]
        \draw (0,0) to[Do] (1,0) to[V, v=$0.65\text{V}$] (2,0) to[R, l=$r_D$] (3.5,0);
    \end{circuitikz}\\
    \vspace{0.2cm}
    \begin{tikzpicture}[scale=0.6]
        \draw[->] (-0.5,0) -- (2.5,0) node[right] {$V_D$};
        \draw[->] (0,-0.5) -- (0,2.5) node[above] {$I_D$};
        \draw[ultra thick, blue] (-0.5,0) -- (0.8,0) -- (2,2);
        \node[below] at (0.8,0) {$0.65\text{V}$};
        \node[right] at (1.4,1) {$m=1/r_D$};
    \end{tikzpicture}\\
    \vspace{0.2cm}
    \raggedright\footnotesize
    \textbf{Directa:} $V_D = V_{D0} + I_D \cdot r_D$.\\
    \textbf{Uso:} Circuitos de potencia donde la corriente es elevada y la caída supera \qty{0.8}{\volt}-\qty{1.0}{\volt}.
\end{minipage}

\vspace{0.5cm}

% ==========================================
% SECCIÓN 2: SEÑALES MIXTAS
% ==========================================
\section{Señales Mixtas y Modelo de Pequeña Señal}
En instrumentación (mecatrónica y biomédica), los circuitos operan con una polarización continua sobre la cual viaja una señal variable. Usamos la \textbf{Notación Estándar IEEE}:
\begin{itemize}
    \item \textbf{DC (Mayúscula-Mayúscula):} $V_D, I_D$ (Valores totales en DC).
    \item \textbf{CA (Minúscula-Minúscula):} $v_d(t), i_d(t)$ (Señales instantáneas alternadas).
    \item \textbf{Total (Minúscula-Mayúscula):} $v_D(t) = V_{DQ} + v_d(t)$ (Superposición total instantánea).
\end{itemize}

Sustituimos la tensión total en la Ecuación de Shockley y aplicamos Serie de Taylor ($e^x = 1 + x + \frac{x^2}{2!} + \dots$):
\begin{equation}
    i_D(t) = I_{DQ} \left[ 1 + \frac{v_d(t)}{\eta V_T} + \frac{1}{2}\left(\frac{v_d(t)}{\eta V_T}\right)^2 + \dots \right] = I_{DQ} + \underbrace{\frac{I_{DQ}}{\eta V_T} v_d(t)}_{\text{Señal lineal } i_d(t)} + \text{Distorsión}
\end{equation}

\textbf{Condición de Linealidad:} El término cuadrático debe ser despreciable frente al lineal ($\frac{v_d(t)}{\eta V_T} \ll 1$). Para tener $<5\%$ de distorsión armónica, exigimos $v_d \le \mathbf{\qtyrange{5}{10}{\milli\volt}}$. Así, el diodo se modela como $r_d = \frac{\eta V_T}{I_{DQ}}$.

% ==========================================
% SECCIÓN 3: PROBLEMA Y METODOLOGÍA
% ==========================================
\section{Análisis Desacoplado: Problema de Ingeniería}
\textbf{Enunciado:} Fuente $V_{CC} = \qty{10}{\volt}$ con ruido $v_s(t) = 0.5 \sin(2\pi \cdot 100 t)\,\text{V}$. $R_S = \qty{1}{\kilo\ohm}$ y Diodo 1N4148 ($\eta = 1.75, V_T = \qty{25.85}{\milli\volt}$). Calcule el punto Q, $r_d$ y la atenuación de rizado.

\vspace{0.2cm}
\noindent
\begin{minipage}{0.45\textwidth}
    \centering \textbf{1. DC (Apagar CA $\implies v_s = 0$)}\\
    \begin{circuitikz}[american, scale=0.7]
        \draw (0,0) to[V, l=$V_{CC}$] (0,2) to[R, l=$R_S$] (3,2) to[D, l=$V_\gamma$, v=$V_{DQ}$] (3,0) -- (0,0);
    \end{circuitikz}\\
    $I_{DQ} = \frac{10 - 0.7}{1000} = \qty{9.3}{\milli\ampere}$\\
    $r_d = \frac{45.24\text{ mV}}{9.3\text{ mA}} = \qty{4.86}{\ohm}$
\end{minipage}
\hfill
\begin{minipage}{0.45\textwidth}
    \centering \textbf{2. CA (Apagar DC $\implies V_{CC} = \text{GND}$)}\\
    \begin{circuitikz}[american, scale=0.7]
        \draw (0,0) to[sV, l=$v_s(t)$] (0,2) to[R, l=$R_S$] (3,2) to[R, l=$r_d$, v=$v_o(t)$] (3,0) -- (0,0);
    \end{circuitikz}\\
    $v_o(t) = 500\text{mV} \cdot \left(\frac{4.86}{1004.86}\right) = \mathbf{2.42\text{ mV}_\text{pk}}$\\
    $\text{Atenuación} = 20\log_{10}\left(\frac{2.42}{500}\right) = \mathbf{\qty{-46.3}{\decibel}}$
\end{minipage}

\begin{aviso}[Conexión PFI (Hitos 1 y 2)]
En las líneas de sensado de corriente (Shunt) y protección ADC, la señal útil viaja sobre un DC. Si polarizamos a alta $I_{DQ}$, $r_d$ cae a pocos ohmios, \textbf{derivando a tierra todo el ruido de conmutación} sin alterar la medición.
\end{aviso}

\newpage
% ==========================================
% SECCIÓN 4: PYTHON Y CONTROL DE IA
% ==========================================
\section{Taller Computacional: Distorsión vs. Linealidad}
Este script evalúa numéricamente el error entre la aproximación lineal de pequeña señal y la resolución exacta no lineal de Shockley, visualizando la Distorsión Armónica Total (THD) si violamos la condición de linealidad.

\begin{lstlisting}
# ==============================================================================
# UNIVERSIDAD CATÓLICA BOLIVIANA "SAN PABLO" - SEDE TARIJA
# Circuitos Electrónicos III (IMT-221) - Semestre 2-2026
# Sesión 4: Validación del Modelo de Pequeña Señal vs. Modelo No Lineal de Shockley
# Docente: M.Sc. Ing. Bernardo Quiroga Turdera
# ==============================================================================
import numpy as np
import matplotlib.pyplot as plt
from scipy.optimize import fsolve

# ------------------------------------------------------------------------------
# 1. PARÁMETROS DEL CIRCUITO Y DEL SEMICONDUCTOR (1N4148)
# ------------------------------------------------------------------------------
Vcc = 10.0           # Nivel DC de alimentación (Volts)
Rs = 1000.0          # Resistencia en serie (Ohms)
Vs_ac_amp = 0.5      # Amplitud pico de la señal de CA (Volts)
frecuencia = 100.0   # Frecuencia de la señal (Hz)

q = 1.602e-19
kB = 1.3806e-23
T_K = 27 + 273.15
Vt = (kB * T_K) / q  # 25.85 mV
Is = 2.52e-9         # Corriente de saturación inversa
eta = 1.75           # Factor de idealidad

# ------------------------------------------------------------------------------
# 2. ANÁLISIS EN DC Y CÁLCULO DE r_d (PEQUEÑA SEÑAL)
# ------------------------------------------------------------------------------
def eq_dc(vd):
    return Is * (np.exp(vd / (eta * Vt)) - 1) - (Vcc - vd) / Rs

Vdq = fsolve(eq_dc, 0.7)[0]
Idq = (Vcc - Vdq) / Rs
rd = (eta * Vt) / Idq

print("==========================================================")
print("             ANÁLISIS DE POLARIZACIÓN EN DC               ")
print("==========================================================")
print(f"Punto Q exacto        : Vdq = {Vdq:.4f} V, Idq = {Idq*1e3:.3f} mA")
print(f"Resistencia dinámica  : rd = {rd:.3f} Ohms")
print(f"Límite de pequeña señal (5% THD): vd_max = {0.10 * eta * Vt * 1e3:.2f} mV")
print("==========================================================")

# ------------------------------------------------------------------------------
# 3. COMPARACIÓN TEMPORAL: EXACTA (NO LINEAL) vs. LINEALIZADA (PEQUEÑA SEÑAL)
# ------------------------------------------------------------------------------
# Vector temporal para 2 ciclos de la señal de 100 Hz
t = np.linspace(0, 2 / frecuencia, 1000)
vs_t = Vs_ac_amp * np.sin(2 * np.pi * frecuencia * t)
vs_total = Vcc + vs_t

# Solución A: Modelo Lineal de Pequeña Señal (Superposición)
vo_pequena_senal = vs_t * (rd / (Rs + rd))
vO_lineal = Vdq + vo_pequena_senal

# Solución B: Solución Exacta Punto a Punto mediante Shockley
vO_exacta = np.zeros_like(t)

for i in range(len(t)):
    def eq_inst(v):
        return Is * (np.exp(v / (eta * Vt)) - 1) - (vs_total[i] - v) / Rs
    vO_exacta[i] = fsolve(eq_inst, Vdq)[0]

# ------------------------------------------------------------------------------
# 4. GRAFICACIÓN Y ANÁLISIS DE ERROR
# ------------------------------------------------------------------------------
plt.figure(figsize=(11, 7))

# Subplot 1: Formas de onda superpuestas
plt.subplot(2, 1, 1)
plt.plot(t * 1e3, (vO_exacta - Vdq) * 1e3, 'b-', linewidth=2, label='Respuesta Exacta No Lineal (Shockley)')
plt.plot(t * 1e3, vo_pequena_senal * 1e3, 'r--', linewidth=2, label='Modelo Lineal de Pequeña Señal ($r_d$)')
plt.title('Comparación de la Señal de CA en el Diodo ($v_d(t)$)', fontsize=12)
plt.ylabel('Amplitud de CA (mV)', fontsize=11)
plt.grid(True, linestyle='--', alpha=0.7)
plt.legend(fontsize=10)

# Subplot 2: Error absoluto de linealización
error_mv = (vO_exacta - vO_lineal) * 1e3
plt.subplot(2, 1, 2)
plt.plot(t * 1e3, error_mv, 'k-', linewidth=1.5, label='Error de Linealización ($v_{exacta} - v_{lineal}$)')
plt.title('Error Absoluto del Modelo Lineal', fontsize=12)
plt.xlabel('Tiempo (ms)', fontsize=11)
plt.ylabel('Error (mV)', fontsize=11)
plt.grid(True, linestyle='--', alpha=0.7)
plt.legend(fontsize=10)

plt.tight_layout()
plt.show()

max_error = np.max(np.abs(error_mv))
print(f"Error máximo absoluto de la aproximación de pequeña señal: {max_error:.4f} mV")
\end{lstlisting}

\begin{aviso}[Protocolo de Auditoría Dinámica (IA)]
En la mesa de trabajo, el docente pedirá modificar la línea 17: \texttt{Vs\_ac\_amp} de \qty{0.5}{\volt} a \qty{8.0}{\volt}.\\
\textbf{Pregunta de Control:} \textit{¿Por qué la forma de onda del error en el Subplot 2 deja de ser simétrica y se deforma?} \\
\textbf{Respuesta Esperada:} Al superar el límite matemático de $\approx \qty{5}{\milli\volt}$ en los bornes del diodo, los términos de segundo y tercer orden de la Serie de Taylor (distorsión armónica) dominan el comportamiento y rompen la asunción de linealidad.
\end{aviso}

\end{document}